\documentclass[11pt]{article}
\usepackage[margin=1in]{geometry}
\usepackage{amsmath, amssymb, amsthm}
\usepackage{color}
\usepackage{hyperref}

\title{Homework 2.3: First-Order Linear Differential Equations}
\author{Tashfeen Omran}
\date{September 2026}

\begin{document}

\maketitle

\tableofcontents
\newpage

\section{Part 1: Comprehensive Introduction, Context, and Prerequisites}

\subsection{Core Concepts}
A differential equation is an equation that relates an unknown function to its derivatives. In this chapter, we focus on \textbf{First-Order Linear Differential Equations}. The term ``first-order'' means that the highest derivative present in the equation is the first derivative (e.g., $\frac{dy}{dx}$ or $y'$). The term ``linear'' means that the dependent variable (typically $y$) and its derivative appear only to the first power and are not multiplied together, nor are they inside non-linear functions like $\sin(y)$ or $e^y$. 

The standard form of a first-order linear differential equation is:
\[ \frac{dy}{dx} + P(x)y = Q(x) \]
where $P(x)$ and $Q(x)$ are continuous functions of the independent variable $x$. If an equation can be algebraically manipulated into this exact format, it is a linear equation.

To solve this type of equation, we introduce a brilliant mathematical concept known as an \textbf{Integrating Factor}, usually denoted by $\mu(x)$ (mu of x). The goal of the integrating factor is to multiply the entire differential equation by a carefully chosen function so that the left side of the equation collapses into the exact derivative of a product, allowing us to simply integrate both sides to find the solution.

\subsection{Intuition and Derivation}
Why does the integrating factor method work? Let us look closely at the Product Rule from differential calculus:
\[ \frac{d}{dx} [\mu(x) y] = \mu(x) \frac{dy}{dx} + \mu'(x) y \]
If we take our standard linear equation, $\frac{dy}{dx} + P(x)y = Q(x)$, and multiply every term by our unknown integrating factor $\mu(x)$, we get:
\[ \mu(x) \frac{dy}{dx} + \mu(x) P(x)y = \mu(x) Q(x) \]
Compare the left side of this equation to the expanded Product Rule. For the left side to be exactly the derivative of the product $\mu(x)y$, we must force the following condition to be true:
\[ \mu'(x) = \mu(x) P(x) \]
This is a separable differential equation for $\mu(x)$:
\[ \frac{d\mu}{dx} = \mu P(x) \implies \frac{d\mu}{\mu} = P(x) dx \]
Integrating both sides yields:
\[ \ln|\mu| = \int P(x) dx \implies \mu(x) = e^{\int P(x) dx} \]
By choosing this specific $\mu(x)$, the differential equation magically simplifies to:
\[ \frac{d}{dx} [\mu(x) y] = \mu(x) Q(x) \]
We can now integrate both sides with respect to $x$ and isolate $y$ to find the general solution.

\subsection{Historical Context and Motivation}
The study of differential equations began in the late 17th century with the founders of calculus, Isaac Newton and Gottfried Wilhelm Leibniz. As scientists attempted to model the physical world—such as the cooling of a hot object, the mixing of chemical solutions, or the mechanics of falling bodies subject to air resistance—they realized that the rate of change of a quantity often depended linearly on the quantity itself. The integrating factor technique was formalized by the Swiss mathematician Leonhard Euler in the 18th century. Euler's method provided a systematic, algebraic algorithm to solve these ubiquitous linear rate-of-change models, fundamentally bridging abstract calculus and applied physics.

\subsection{Key Formulas}
\begin{enumerate}
    \item \textbf{Standard Form}: $\frac{dy}{dx} + P(x)y = Q(x)$
    \item \textbf{Integrating Factor}: $\mu(x) = e^{\int P(x) dx}$
    \item \textbf{Collapsed Form}: $\frac{d}{dx} [\mu(x) y] = \mu(x) Q(x)$
    \item \textbf{General Solution}: $y(x) = \frac{1}{\mu(x)} \left( \int \mu(x) Q(x) dx + C \right)$
\end{enumerate}

\subsection{Prerequisites}
To succeed in this topic, you must have a firm grasp of:
\begin{itemize}
    \item \textbf{Algebraic manipulation}: Distributing terms, factoring, and working with rational expressions.
    \item \textbf{Properties of Exponents and Logarithms}: Crucially, $e^{\ln(x)} = x$ and $e^{a \ln(x)} = x^a$.
    \item \textbf{Integration Techniques}: Basic antiderivatives, $u$-substitution, and integration by parts (often required when integrating $\mu(x) Q(x)$).
    \item \textbf{The Product Rule of Differentiation}: Recognizing the expanded form of a product derivative is essential for understanding the setup.
\end{itemize}

\newpage
\section{Part 2: Detailed Homework Solutions}

\subsection*{Problem 1}
\textbf{Question:} Determine whether the given equation is separable, linear, neither, or both.
\[ \frac{dx}{dt} - 6xt = 4e^x \]
Choose the correct answer below.

\textbf{Solution:}
Let us analyze the dependent variable $x$ and the independent variable $t$.
1. \textbf{Linearity check}: For the equation to be linear in $x$, all instances of $x$ and its derivatives must be to the first power and cannot be arguments of non-linear functions. However, the term on the right side is $4e^x$. Because $x$ is trapped in the exponent of $e$, the equation is \textbf{non-linear}.
2. \textbf{Separability check}: An equation is separable if it can be written as $\frac{dx}{dt} = g(t)h(x)$. Let's isolate the derivative:
\[ \frac{dx}{dt} = 6xt + 4e^x \]
We cannot factor the expression $6xt + 4e^x$ into a product of a function of purely $t$ and a function of purely $x$. Therefore, the equation is \textbf{not separable}.

\textbf{Final Answer:} \textbf{C. The equation is neither separable nor linear.}

\hrulefill

\subsection*{Problem 2}
\textbf{Question:} Determine whether the given equation is separable, linear, neither, or both.
\[ 2t = 4e^t \frac{dy}{dt} + 5y \ln t \]
Choose the correct answer below.

\textbf{Solution:}
Let us analyze the dependent variable $y$ and the independent variable $t$.
1. \textbf{Linearity check}: We look at the terms involving $y$ and $\frac{dy}{dt}$. They appear strictly to the first power. The non-linear functions ($e^t$ and $\ln t$) only apply to the independent variable $t$, which is perfectly acceptable for a linear equation. Let's rewrite it in standard form to be sure:
\[ 4e^t \frac{dy}{dt} + (5\ln t)y = 2t \implies \frac{dy}{dt} + \left( \frac{5\ln t}{4e^t} \right)y = \frac{2t}{4e^t} \]
This matches the standard form $\frac{dy}{dt} + P(t)y = Q(t)$. Thus, it is \textbf{linear}.
2. \textbf{Separability check}: Let's isolate the derivative:
\[ 4e^t \frac{dy}{dt} = 2t - 5y \ln t \implies \frac{dy}{dt} = \frac{2t - 5y \ln t}{4e^t} \]
The numerator cannot be factored into a purely $y$ component and a purely $t$ component. Thus, it is \textbf{not separable}.

\textbf{Final Answer:} \textbf{A. The equation is linear.}

\hrulefill

\subsection*{Problem 3}
\textbf{Question:} Determine whether the given equation is separable, linear, neither, or both.
\[ 2x \frac{dx}{dt} + 5t^3 x = 8\sin t \]
Choose the correct answer below.

\textbf{Solution:}
Let us analyze the dependent variable $x$ and the independent variable $t$.
1. \textbf{Linearity check}: Look at the first term: $2x \frac{dx}{dt}$. Here, the dependent variable $x$ is multiplied by its own derivative. This violates the rule for linearity. Thus, the equation is \textbf{non-linear}.
2. \textbf{Separability check}: Let us attempt to isolate $\frac{dx}{dt}$:
\[ 2x \frac{dx}{dt} = 8\sin t - 5t^3 x \implies \frac{dx}{dt} = \frac{8\sin t - 5t^3 x}{2x} = \frac{4\sin t}{x} - \frac{5}{2}t^3 \]
This expression cannot be factored into $g(t)h(x)$. It is \textbf{not separable}.

\textbf{Final Answer:} \textbf{C. The equation is neither separable nor linear.}

\hrulefill

\subsection*{Problem 4}
\textbf{Question:} Determine whether the given equation is separable, linear, neither, or both.
\[ r = \frac{dr}{d\theta} - 7\theta^3 \]
Choose the correct answer below.

\textbf{Solution:}
Let us analyze the dependent variable $r$ and independent variable $\theta$.
1. \textbf{Linearity check}: Rearrange the equation to group $r$ terms:
\[ \frac{dr}{d\theta} - r = 7\theta^3 \]
This is in the exact standard form $\frac{dr}{d\theta} + P(\theta)r = Q(\theta)$, where $P(\theta) = -1$ and $Q(\theta) = 7\theta^3$. It is \textbf{linear}.
2. \textbf{Separability check}: Isolate the derivative:
\[ \frac{dr}{d\theta} = r + 7\theta^3 \]
The addition makes it impossible to factor this into a function of $r$ times a function of $\theta$. It is \textbf{not separable}.

\textbf{Final Answer:} \textbf{A. The equation is linear.}

\hrulefill

\subsection*{Problem 5}
\textbf{Question:} Obtain the general solution to the equation.
\[ \frac{dy}{dx} = \frac{y}{x} + 9x + 8 \]

\textbf{Solution:}
\textbf{Step 1: Put the equation in standard form.}
Subtract $\frac{y}{x}$ from both sides:
\[ \frac{dy}{dx} - \frac{1}{x}y = 9x + 8 \]
Now it matches $\frac{dy}{dx} + P(x)y = Q(x)$, where $P(x) = -\frac{1}{x}$ and $Q(x) = 9x + 8$.

\textbf{Step 2: Find the Integrating Factor $\mu(x)$.}
\[ \mu(x) = e^{\int P(x) dx} = e^{\int -\frac{1}{x} dx} = e^{-\ln(x)} = e^{\ln(x^{-1})} = x^{-1} = \frac{1}{x} \]

\textbf{Step 3: Multiply the standard form equation by $\mu(x)$.}
\[ \frac{1}{x} \frac{dy}{dx} - \frac{1}{x^2} y = \frac{1}{x} (9x + 8) \]
The left side collapses to the derivative of $(\mu(x) y)$:
\[ \frac{d}{dx} \left[ x^{-1} y \right] = 9 + \frac{8}{x} \]

\textbf{Step 4: Integrate both sides with respect to $x$.}
\[ \int \frac{d}{dx} \left[ x^{-1} y \right] dx = \int \left( 9 + \frac{8}{x} \right) dx \]
\[ x^{-1} y = 9x + 8\ln|x| + C \]

\textbf{Step 5: Solve for $y$.}
Multiply both sides by $x$:
\[ y(x) = x(9x + 8\ln|x| + C) \]
\[ y(x) = 9x^2 + 8x\ln|x| + Cx \]

\textbf{Final Answer:} $y(x) = 9x^2 + 8x\ln|x| + Cx$

\hrulefill

\subsection*{Problem 6}
\textbf{Question:} Obtain the general solution to the equation.
\[ (5t + y + 1)dt - dy = 0 \]

\textbf{Solution:}
\textbf{Step 1: Put the equation in standard form.}
First, divide everything by $dt$ to reveal the derivative:
\[ 5t + y + 1 - \frac{dy}{dt} = 0 \implies \frac{dy}{dt} - y = 5t + 1 \]
Here, $P(t) = -1$ and $Q(t) = 5t + 1$.

\textbf{Step 2: Find the Integrating Factor $\mu(t)$.}
\[ \mu(t) = e^{\int -1 dt} = e^{-t} \]

\textbf{Step 3: Multiply the standard form equation by $\mu(t)$.}
\[ e^{-t} \frac{dy}{dt} - e^{-t} y = e^{-t} (5t + 1) \]
Collapse the left side:
\[ \frac{d}{dt} \left[ e^{-t} y \right] = (5t + 1)e^{-t} \]

\textbf{Step 4: Integrate both sides.}
\[ e^{-t} y = \int (5t + 1)e^{-t} dt \]
We must use \textbf{Integration by Parts} for the right side: $\int u \, dv = uv - \int v \, du$.
Let $u = 5t + 1 \implies du = 5 \, dt$
Let $dv = e^{-t} dt \implies v = -e^{-t}$
\[ \int (5t + 1)e^{-t} dt = (5t + 1)(-e^{-t}) - \int (-e^{-t})(5) dt \]
\[ = -(5t + 1)e^{-t} + 5 \int e^{-t} dt \]
\[ = -5te^{-t} - e^{-t} - 5e^{-t} + C = -5te^{-t} - 6e^{-t} + C \]
So we have:
\[ e^{-t} y = -e^{-t}(5t + 6) + C \]

\textbf{Step 5: Solve for $y$.}
Multiply both sides by $e^t$:
\[ y(t) = - (5t + 6) + Ce^t = -5t - 6 + Ce^t \]

\textbf{Final Answer:} $y(t) = -5t - 6 + Ce^t$

\hrulefill

\subsection*{Problem 7}
\textbf{Question:} Obtain the general solution to the equation.
\[ y \frac{dx}{dy} + 8x = 7y^5 \]

\textbf{Solution:}
\textbf{Step 1: Put the equation in standard form.}
Notice that $x$ is the dependent variable and $y$ is the independent variable. Divide by $y$ to make the leading coefficient 1:
\[ \frac{dx}{dy} + \frac{8}{y} x = 7y^4 \]
This is linear in $x$ with $P(y) = \frac{8}{y}$ and $Q(y) = 7y^4$.

\textbf{Step 2: Find the Integrating Factor $\mu(y)$.}
\[ \mu(y) = e^{\int \frac{8}{y} dy} = e^{8\ln|y|} = e^{\ln(y^8)} = y^8 \]

\textbf{Step 3: Multiply the standard form equation by $\mu(y)$.}
\[ y^8 \frac{dx}{dy} + 8y^7 x = 7y^{12} \]
Collapse the left side:
\[ \frac{d}{dy} \left[ y^8 x \right] = 7y^{12} \]

\textbf{Step 4: Integrate both sides with respect to $y$.}
\[ \int \frac{d}{dy} \left[ y^8 x \right] dy = \int 7y^{12} dy \]
\[ y^8 x = \frac{7}{13}y^{13} + C \]

\textbf{Step 5: Solve for $x$.}
Divide by $y^8$:
\[ x(y) = \frac{7}{13}y^5 + Cy^{-8} \]

\textbf{Final Answer:} $x(y) = \frac{7}{13}y^5 + Cy^{-8}$

\hrulefill

\subsection*{Problem 8}
\textbf{Question:} Solve the initial value problem.
\[ \frac{dy}{dx} - \frac{y}{x} = 8xe^x, \quad y(1) = 8e - 6 \]

\textbf{Solution:}
\textbf{Step 1: Identify standard form.}
The equation is already in standard form: $\frac{dy}{dx} - \frac{1}{x}y = 8xe^x$.
Here, $P(x) = -\frac{1}{x}$ and $Q(x) = 8xe^x$.

\textbf{Step 2: Find the Integrating Factor $\mu(x)$.}
\[ \mu(x) = e^{\int -\frac{1}{x} dx} = e^{-\ln(x)} = \frac{1}{x} \]

\textbf{Step 3: Multiply the standard form equation by $\mu(x)$.}
\[ \frac{1}{x} \frac{dy}{dx} - \frac{1}{x^2}y = \frac{1}{x} (8xe^x) \]
\[ \frac{d}{dx} \left[ x^{-1} y \right] = 8e^x \]

\textbf{Step 4: Integrate both sides.}
\[ x^{-1} y = \int 8e^x dx \]
\[ x^{-1} y = 8e^x + C \]

\textbf{Step 5: Solve for $y$ (General Solution).}
Multiply by $x$:
\[ y(x) = 8xe^x + Cx \]

\textbf{Step 6: Apply the Initial Condition $y(1) = 8e - 6$.}
Substitute $x = 1$ and $y = 8e - 6$:
\[ 8e - 6 = 8(1)e^1 + C(1) \]
\[ 8e - 6 = 8e + C \]
Subtract $8e$ from both sides:
\[ C = -6 \]

\textbf{Step 7: Write the final particular solution.}
\[ y(x) = 8xe^x - 6x \]

\textbf{Final Answer:} $y(x) = 8xe^x - 6x$

\hrulefill

\subsection*{Problem 9}
\textbf{Question:} Solve the initial value problem.
\[ \frac{dy}{dx} + 5y - 5e^{-x} = 0, \quad y(0) = \frac{9}{4} \]

\textbf{Solution:}
\textbf{Step 1: Put the equation in standard form.}
Move $-5e^{-x}$ to the right side:
\[ \frac{dy}{dx} + 5y = 5e^{-x} \]
Here, $P(x) = 5$ and $Q(x) = 5e^{-x}$.

\textbf{Step 2: Find the Integrating Factor $\mu(x)$.}
\[ \mu(x) = e^{\int 5 dx} = e^{5x} \]

\textbf{Step 3: Multiply the standard form equation by $\mu(x)$.}
\[ e^{5x} \frac{dy}{dx} + 5e^{5x} y = e^{5x}(5e^{-x}) \]
Collapse the left side and simplify the right side using exponent rules ($e^{5x} \cdot e^{-x} = e^{4x}$):
\[ \frac{d}{dx} \left[ e^{5x} y \right] = 5e^{4x} \]

\textbf{Step 4: Integrate both sides.}
\[ e^{5x} y = \int 5e^{4x} dx = \frac{5}{4}e^{4x} + C \]

\textbf{Step 5: Solve for $y$ (General Solution).}
Divide both sides by $e^{5x}$ (or multiply by $e^{-5x}$):
\[ y(x) = \frac{5}{4}e^{4x-5x} + Ce^{-5x} = \frac{5}{4}e^{-x} + Ce^{-5x} \]

\textbf{Step 6: Apply the Initial Condition $y(0) = \frac{9}{4}$.}
Substitute $x = 0$ and $y = \frac{9}{4}$:
\[ \frac{9}{4} = \frac{5}{4}e^0 + Ce^0 \]
\[ \frac{9}{4} = \frac{5}{4} + C \implies C = \frac{9}{4} - \frac{5}{4} = \frac{4}{4} = 1 \]

\textbf{Step 7: Write the final particular solution.}
\[ y(x) = \frac{5}{4}e^{-x} + e^{-5x} \]

\textbf{Final Answer:} $y(x) = \frac{5}{4}e^{-x} + e^{-5x}$

\hrulefill

\subsection*{Problem 10}
\textbf{Question:} Solve the initial value problem.
\[ t^2 \frac{dx}{dt} + 2tx = t^4 \ln t + 2, \quad x(1) = 0 \]

\textbf{Solution:}
\textbf{Step 1: Analyze the given equation.}
Take a close look at the left side of the equation: $t^2 \frac{dx}{dt} + 2tx$. 
Notice that this is exactly the expanded form of the Product Rule for the product $(t^2 \cdot x)$. 
\[ \frac{d}{dt} \left[ t^2 x \right] = t^2 \frac{dx}{dt} + 2t x \]
Because the equation is already perfectly set up, we do not need to divide by $t^2$ to find an integrating factor; the equation is already multiplied by it! We can immediately collapse the left side.

\textbf{Step 2: Collapse and Integrate.}
\[ \frac{d}{dt} \left[ t^2 x \right] = t^4 \ln t + 2 \]
Integrate both sides with respect to $t$:
\[ t^2 x = \int (t^4 \ln t + 2) dt = \int t^4 \ln t \, dt + \int 2 \, dt \]
The integral of $2$ is just $2t$. For $\int t^4 \ln t \, dt$, we use \textbf{Integration by Parts}:
Let $u = \ln t \implies du = \frac{1}{t} dt$
Let $dv = t^4 dt \implies v = \frac{1}{5}t^5$
\[ \int t^4 \ln t \, dt = (\ln t)\left(\frac{1}{5}t^5\right) - \int \left(\frac{1}{5}t^5\right)\left(\frac{1}{t}\right) dt \]
\[ = \frac{1}{5}t^5 \ln t - \frac{1}{5} \int t^4 dt = \frac{1}{5}t^5 \ln t - \frac{1}{25}t^5 \]
Putting it all together:
\[ t^2 x = \frac{1}{5}t^5 \ln t - \frac{1}{25}t^5 + 2t + C \]

\textbf{Step 3: Solve for $x(t)$.}
Divide the entire equation by $t^2$:
\[ x(t) = \frac{1}{5}t^3 \ln t - \frac{1}{25}t^3 + \frac{2}{t} + \frac{C}{t^2} \]

\textbf{Step 4: Apply Initial Condition $x(1) = 0$.}
Substitute $t = 1$ and $x = 0$:
\[ 0 = \frac{1}{5}(1)^3 \ln(1) - \frac{1}{25}(1)^3 + \frac{2}{1} + \frac{C}{(1)^2} \]
Since $\ln(1) = 0$, the first term vanishes:
\[ 0 = 0 - \frac{1}{25} + 2 + C \]
\[ 0 = \frac{49}{25} + C \implies C = -\frac{49}{25} \]

\textbf{Step 5: Write final particular solution.}
\[ x(t) = \frac{1}{5}t^3 \ln t - \frac{1}{25}t^3 + \frac{2}{t} - \frac{49}{25 t^2} \]

\textbf{Final Answer:} $x(t) = \frac{1}{5}t^3 \ln t - \frac{1}{25}t^3 + 2t^{-1} - \frac{49}{25}t^{-2}$

\newpage
\section{Part 3: In-Depth Analysis of Problems and Techniques}

\subsection{A) Problem Types and General Approach}
The homework problems map nicely into three distinct categories:
\begin{enumerate}
    \item \textbf{Classification Problems (Q1-Q4):} The objective is simply to diagnose the equation. 
    \textit{Approach:} Identify the dependent variable and the independent variable. Check for linearity by ensuring the dependent variable and its derivative appear only to the first power. Check separability by attempting to algebraically isolate the derivative and factor the remaining side into a product of purely independent and purely dependent functions.
    \item \textbf{General Solution Problems (Q5-Q7):} We must find the infinite family of solutions with an arbitrary constant $C$.
    \textit{Approach:} Force the equation into standard form $\frac{dy}{dx} + P(x)y = Q(x)$ by dividing by the leading coefficient if necessary. Calculate $\mu(x) = e^{\int P(x)dx}$. Multiply through, collapse the left side using the reverse product rule, integrate both sides, and solve algebraically for the dependent variable.
    \item \textbf{Initial Value Problems (IVPs) (Q8-Q10):} We are given a specific coordinate that the function passes through, which eliminates the unknown $+C$. 
    \textit{Approach:} Solve exactly like a General Solution problem up to the very end. Once you have the general equation with $+C$, plug in the initial $x$ and $y$ values, isolate $C$, and substitute it back into the general solution.
\end{enumerate}

\subsection{B) Key Algebraic and Calculus Manipulations}
\begin{itemize}
    \item \textbf{Logarithm and Exponent Cancellation (Q5, Q7, Q8):} A persistent trick is simplifying $\mu(x)$ using the rule $e^{\ln(f(x))} = f(x)$. Furthermore, we used the power rule for logarithms $a\ln(x) = \ln(x^a)$ in Q7: $e^{8\ln y} \to e^{\ln(y^8)} \to y^8$. Failure to recognize this trick will leave you attempting to integrate an impossible expression.
    \item \textbf{Integration by Parts (Q6, Q10):} Integrating the right-hand side after multiplying by $\mu(x)$ frequently results in products of dissimilar functions, such as an exponential multiplied by a polynomial (Q6) or a logarithm multiplied by a polynomial (Q10). The formula $\int u \, dv = uv - \int v \, du$ is absolutely essential.
    \item \textbf{Variable Reversal (Q7):} Question 7 elegantly flips the script by making $x$ the dependent variable and $y$ the independent variable. The calculus logic is identical, but tracking $dx/dy$ and ensuring $P$ and $Q$ are strictly functions of $y$ prevents confusion.
    \item \textbf{Pre-Collapsed Product Rule (Q10):} Question 10 is uniquely designed to test conceptual understanding over raw algorithmic execution. By recognizing that $t^2 \frac{dx}{dt} + 2tx$ is already the exact derivative $\frac{d}{dt}[t^2 x]$, we completely bypass the need to define $P(t)$, find an integrating factor, and multiply the equation. 
\end{itemize}

\section{Part 4: "Cheatsheet" and Tips for Success}

\subsection{Formulas at a Glance}
\begin{itemize}
    \item \textbf{Standard Linear DE:} $y' + P(x)y = Q(x)$
    \item \textbf{Integrating Factor:} $\mu(x) = e^{\int P(x) dx}$
    \item \textbf{Collapsed Form:} $(\mu y)' = \mu Q$
\end{itemize}

\subsection{Pro-Tips and Tricks}
\begin{itemize}
    \item \textbf{The Negative Sign Cheat:} When calculating $\mu(x)$, remember that the sign of $P(x)$ must be included. If you have $y' - 5y = x$, then $P(x) = -5$, not $5$.
    \item \textbf{The Absolute Value Shortcut:} When evaluating $\mu(x) = e^{\int \frac{1}{x} dx} = e^{\ln|x|}$, it is standard practice to assume $x > 0$ yielding just $x$, unless initial conditions dictate otherwise.
    \item \textbf{The "Plus C" Placement Rule:} Always add the $+C$ \textit{immediately} after you perform the integration step, while the integrating factor is still attached to the dependent variable. (i.e. $\mu y = \int \dots + C$). NEVER wait until the end to slap a $+C$ on the final answer, because the $C$ must also be divided by $\mu(x)$ when isolating $y$.
\end{itemize}

\subsection{Common Pitfalls}
\begin{itemize}
    \item \textbf{The Leading Coefficient Trap:} (e.g., Q7) Do not calculate $\mu(x)$ if the derivative term has a coefficient. You \textit{must} divide the entire equation by the coefficient first so that the derivative stands alone.
    \item \textbf{Algebraic Sign Errors when Moving Terms:} (e.g., Q5, Q9) Frequently, the equation starts with the $P(x)y$ term on the wrong side. Subtracting it over to the left side flips its sign. Students often miss this and compute a $\mu(x)$ with the wrong sign.
\end{itemize}

\newpage
\section{Part 5: Conceptual Synthesis and The "Big Picture"}

\subsection{A) Thematic Connections}
The core theme of the integrating factor technique is \textbf{mathematical transformation to reveal structure}. In early algebra, you multiplied both sides of an equation by a common denominator to clear fractions. In integral calculus, you use $u$-substitution to compress a complex integrand into a simple one. Here, we multiply a messy differential equation by a specifically engineered function ($\mu$) designed to artificially construct a perfect Product Rule derivative. We are transforming a problem we do not know how to solve into one that can be solved by simply taking an anti-derivative.

\subsection{B) Forward and Backward Links}
\textbf{Backward:} This topic is the direct climax of differential calculus (Product Rule) combined with integral calculus (Integration by Parts, Exponential rules). It completely unites both halves of Calculus I and II into an algebraic procedure.

\textbf{Forward:} The conceptual leap of "finding a multiplier to force a perfect derivative" is exactly the foundation for the study of \textbf{Exact Equations} in multivariable calculus. Furthermore, finding standard forms and generating specific transient vs. steady-state solutions (seen in Q9 where $e^{-5x}$ decays to 0) directly precedes \textbf{Laplace Transforms} and the study of dynamic systems and controls in engineering.

\section{Part 6: Real-World Application and Modeling}

\subsection{A) Concrete Scenarios in Finance and Economics}
\begin{enumerate}
    \item \textbf{Continuous Annuity / Cash Flow Models:} In corporate finance, when analyzing continuous income streams, a bank account accrues interest at a continuous rate, and cash is simultaneously injected or withdrawn continuously. Determining the value of the portfolio at any time $t$ requires solving a linear first-order ODE.
    \item \textbf{The Vasicek Interest Rate Model (Stochastic Calculus):} In mathematical finance, interest rates do not stay constant; they fluctuate but tend to revert to a long-term mean (mean-reversion). The expected value of the interest rate under the Vasicek model follows a deterministic linear first-order differential equation.
    \item \textbf{Solow-Swan Economic Growth Model:} In macroeconomics, modeling the capital accumulation of an economy over time involves capital depreciation and continuous investment. The fundamental equation describing capital per capita is a linear differential equation.
\end{enumerate}

\subsection{B) Model Problem Setup: Continuous Cash Flow}
Imagine you are setting up a retirement account. The account guarantees a continuously compounded interest rate of $r = 5\%$ per year. Additionally, to combat inflation, you continuously deposit money into the account at a rate that grows over time: $D(t) = 5000e^{0.03t}$ dollars per year. 

\textbf{Variables:}
\begin{itemize}
    \item $t$: Time in years
    \item $V(t)$: Total value of the portfolio at time $t$
    \item $r = 0.05$: Continuous interest rate
    \item $D(t) = 5000e^{0.03t}$: The continuous income stream
\end{itemize}
\textbf{The Model Formulation:}
The rate of change of the portfolio's value, $\frac{dV}{dt}$, is the sum of the interest generated by the current balance ($0.05V$) plus the incoming cash deposits ($D(t)$).
\[ \frac{dV}{dt} = 0.05V + 5000e^{0.03t} \]
To solve this, we arrange it into standard linear form:
\[ \frac{dV}{dt} - 0.05V = 5000e^{0.03t} \]
An actuary would find $\mu(t) = e^{-0.05t}$ to determine the portfolio's value $V(t)$ decades into the future.

\newpage
\section{Part 7: Common Variations and Untested Concepts}

A vital concept taught in standard first-order differential equations that did not appear in your homework is the \textbf{Bernoulli Differential Equation}.

\subsection{Bernoulli Equations}
A Bernoulli equation is a \textit{non-linear} differential equation of a very specific form that can be magically transformed into a linear one via substitution. The standard form is:
\[ \frac{dy}{dx} + P(x)y = Q(x)y^n \]
where $n$ is any real number other than 0 or 1. (If $n=0$ or $n=1$, it is just a standard linear equation). Because of the $y^n$ term on the right, the equation is decidedly non-linear.

\textbf{The Trick:} We use a universal algebraic substitution: let $v = y^{1-n}$. By applying the chain rule, this substitution brilliantly linearizes the differential equation into the form $\frac{dv}{dx} + (1-n)P(x)v = (1-n)Q(x)$, which can then be solved using the standard Integrating Factor method you practiced in this homework!

\textbf{Brief Example:}
Solve $\frac{dy}{dx} + \frac{1}{x}y = xy^2$.
Notice $n = 2$.
Let $v = y^{1-2} = y^{-1}$. Then $\frac{dv}{dx} = -y^{-2} \frac{dy}{dx}$.
Divide the original equation by $y^2$:
\[ y^{-2}\frac{dy}{dx} + \frac{1}{x}y^{-1} = x \]
Substitute $v$ and $\frac{dv}{dx}$:
\[ -\frac{dv}{dx} + \frac{1}{x}v = x \implies \frac{dv}{dx} - \frac{1}{x}v = -x \]
This is now a purely linear equation for $v$! Once $v$ is found via the integrating factor, substitute back $y = v^{-1}$ to find the final answer.

\newpage
\section{Part 8: Advanced Diagnostic Testing: "Find the Flaw"}

Below are 5 advanced problems. The provided solutions look convincing but each contains one fatal, subtle flaw. \textbf{Your task is to identify the error, explain why it is wrong, and provide the correct calculation.}

\subsection*{Problem 1}
\textbf{Solve:} $x \frac{dy}{dx} + 2y = x^3$
\textbf{Flawed Solution:}
1. The equation is linear. We identify $P(x) = 2$ and $Q(x) = x^3$.
2. Calculate the integrating factor: $\mu(x) = e^{\int 2 dx} = e^{2x}$.
3. Multiply both sides by $\mu(x)$: $e^{2x} x y' + 2e^{2x} y = x^3 e^{2x}$.
4. Collapse the left side: $\frac{d}{dx} [x y e^{2x}] = x^3 e^{2x}$. 
\textit{(Error occurs near here.)}

\subsection*{Problem 2}
\textbf{Solve:} $\frac{dy}{dx} - y = e^x$
\textbf{Flawed Solution:}
1. Standard form: $P(x) = -1$, $Q(x) = e^x$.
2. Integrating factor: $\mu(x) = e^{\int -1 dx} = e^{-x}$.
3. Multiply through: $e^{-x} y' - e^{-x} y = e^{-x} e^x \implies (y e^{-x})' = 1$.
4. Integrate: $y e^{-x} = \int 1 dx \implies y e^{-x} = x$.
5. Isolate $y$: $y = x e^x$.
6. Add the constant of integration: $y = x e^x + C$.
\textit{(Error occurs near here.)}

\subsection*{Problem 3}
\textbf{Solve:} $\frac{dy}{dx} + \frac{1}{x}y = 2$
\textbf{Flawed Solution:}
1. Standard form: $P(x) = \frac{1}{x}$, $Q(x) = 2$.
2. Integrating factor: $\mu(x) = e^{\int \frac{1}{x} dx} = e^{\ln(x)} = x$.
3. Multiply through: $x y' + y = 2x$.
4. Collapse: $(xy)' = 2x$.
5. Integrate: $xy = x^2$.
6. Isolate $y$ and add $C$: $y = x + C$.
\textit{(Error occurs near here.)}

\subsection*{Problem 4}
\textbf{Solve:} $\frac{dy}{dx} - \frac{2}{x} y = x$
\textbf{Flawed Solution:}
1. Standard form: $P(x) = \frac{2}{x}$, $Q(x) = x$.
2. Integrating factor: $\mu(x) = e^{\int \frac{2}{x} dx} = e^{2\ln(x)} = x^2$.
3. Multiply through: $x^2 y' - 2x y = x^3$.
4. Collapse: $(x^2 y)' = x^3$.
5. Integrate: $x^2 y = \frac{1}{4}x^4 + C \implies y = \frac{1}{4}x^2 + Cx^{-2}$.
\textit{(Error occurs near here.)}

\subsection*{Problem 5}
\textbf{Solve:} $\cos(x) \frac{dy}{dx} + \sin(x) y = 1$
\textbf{Flawed Solution:}
1. Standard form: Divide by $\cos(x)$. $y' + \tan(x) y = \sec(x)$.
2. Integrating factor: $\mu(x) = e^{\int \tan(x) dx} = e^{\ln|\cos(x)|} = \cos(x)$.
3. Multiply through: $\cos(x)y' + \sin(x)y = 1$.
4. Collapse: $(\cos(x)y)' = 1 \implies \cos(x)y = x + C$.
5. Isolate $y$: $y = x\sec(x) + C\sec(x)$.
\textit{(Error occurs near here.)}

\hrulefill

\subsection*{Answer Key to "Find the Flaw"}
\textbf{Problem 1 Flaw:} The student identified $P(x)$ before putting the equation into standard form (the leading coefficient was $x$, not 1). 
\textbf{Correction:} First divide by $x$: $y' + \frac{2}{x}y = x^2$. $\mu(x) = x^2$. Final answer: $y = \frac{1}{5}x^3 + Cx^{-2}$.

\textbf{Problem 2 Flaw:} The constant $+C$ was added at the very end after the integrating factor $e^{-x}$ was divided out.
\textbf{Correction:} Add $+C$ immediately upon integrating: $y e^{-x} = x + C$. Thus, isolating $y$ yields $y = x e^x + C e^x$.

\textbf{Problem 3 Flaw:} Similar to Problem 2, the $+C$ was omitted during the integration step and appended arbitrarily at the end.
\textbf{Correction:} $\int 2x dx = x^2 + C$. Then dividing by $x$ gives $y = x + \frac{C}{x}$. 

\textbf{Problem 4 Flaw:} The negative sign in front of the $P(x)$ term was dropped when calculating the integrating factor. $P(x) = -\frac{2}{x}$.
\textbf{Correction:} $\mu(x) = e^{\int -\frac{2}{x} dx} = x^{-2}$. Final answer: $y = x^2 \ln|x| + Cx^2$.

\textbf{Problem 5 Flaw:} The integral of $\tan(x)$ is evaluated incorrectly. $\int \tan(x) dx = -\ln|\cos(x)|$ or $\ln|\sec(x)|$, not $\ln|\cos(x)|$.
\textbf{Correction:} $\mu(x) = e^{\ln|\sec(x)|} = \sec(x)$. Then $\frac{d}{dx}[\sec(x)y] = \sec^2(x) \implies \sec(x)y = \tan(x) + C \implies y = \sin(x) + C\cos(x)$.

\end{document}